\documentclass[11pt,a4paper]{article}
\usepackage[T1]{fontenc}
\usepackage[utf8]{inputenc}
\usepackage{lmodern}
\usepackage[margin=25mm]{geometry}
\usepackage{microtype}
\usepackage{booktabs,longtable,array,tabularx}
\usepackage{xcolor}
\usepackage{graphicx}
\usepackage{enumitem}
\usepackage{url}
\usepackage{xurl}
\usepackage[style=numeric,sorting=none,backend=biber]{biblatex}
\usepackage[colorlinks=true,linkcolor=teal,citecolor=teal,urlcolor=teal]{hyperref}
\addbibresource{references.bib}
\setlength{\emergencystretch}{3em}
\setlength{\parskip}{0.35em}
\setlength{\parindent}{0pt}
\setlist{nosep}
\hypersetup{pdftitle={What About the Independents Then: Investigating Small Large Language Models in the UK},pdfauthor={Phillip Charles O'Brien}}
\newcommand{\CensusVersion}{1.2.1}
\newcommand{\CensusRevisionDate}{2 November 2026}
\newcommand{\CensusModels}{917}
\newcommand{\CensusAccounts}{495}
\newcommand{\UKBuilders}{5}
\newcommand{\UKBuiltModels}{7}
\newcommand{\ScratchBuilders}{3}
\newcommand{\ScratchModels}{5}
\newcommand{\StrictBuilders}{1}
\newcommand{\UKInstitutional}{4}
\newcommand{\PrimaryPeople}{19}
\newcommand{\PrimaryInstitutional}{17}
\newcommand{\StrictModels}{1}
\newcommand{\HistoricalContextPeople}{16}
\newcommand{\PrimaryUniversityYes}{15}
\newcommand{\PrimaryUniversityNo}{1}
\newcommand{\PrimaryUniversityUnknown}{3}
\newcommand{\DirectUKBuildModels}{1}
\newcommand{\InferredUKBuildModels}{6}
\newcommand{\DetailedModelAssessments}{55}
\newcommand{\ReviewedPeople}{35}
\newcommand{\ScratchInventoryModels}{457}
\newcommand{\DerivedInventoryModels}{352}
\newcommand{\MirrorInventoryModels}{36}
\newcommand{\UntrainedInventoryModels}{60}
\newcommand{\UnresolvedProvenanceModels}{12}
\newcommand{\InferredProvenanceModels}{237}
\newcommand{\InferredScratchModels}{154}
\newcommand{\UnresolvedTrainingModels}{15}
\newcommand{\EstimatedTrainedModels}{9}
\newcommand{\UnresolvedOwnScratchLineageModels}{147}
\newcommand{\PartialOwnScratchLineageModels}{2}
\newcommand{\InsideAICentreBuilders}{3}
\newcommand{\OutsideAICentreBuilders}{1}
\newcommand{\UnknownBuildCityBuilders}{1}
\newcommand{\UKSchoolingBuilders}{1}
\newcommand{\NonUKSchoolingBuilders}{2}
\newcommand{\UnknownSchoolingBuilders}{2}
\newcommand{\UKSourceConnectionAccounts}{10}
\newcommand{\UKEstimateAccounts}{2}
\newcommand{\UKAffiliationAccounts}{1}
\newcommand{\MixedConnectionAccounts}{11}
\newcommand{\NonUKConnectionOrEstimateAccounts}{387}
\newcommand{\RegionClueAccounts}{9}
\newcommand{\UnresolvedGeographyAccounts}{75}
\newcommand{\UnresolvedGeographyModels}{98}
\newcommand{\FollowthroughAccounts}{151}
\newcommand{\CountryEstimateAccounts}{66}
\newcommand{\HighCountryEstimateAccounts}{23}
\newcommand{\MediumCountryEstimateAccounts}{26}
\newcommand{\LowCountryEstimateAccounts}{17}
\newcommand{\UnknownAccounts}{75}
\newcommand{\UnknownAccountModels}{98}


\begin{document}
% Match the crest-first cover used in the author's COLM and WiggleGPT papers.
\hypersetup{pageanchor=false}
\begin{titlepage}
\setlength{\parskip}{0pt}
\centering
\vspace*{1cm}
\includegraphics[width=0.65\textwidth]{figures/obrien-crest.png}
\vspace{1cm}

{\LARGE\bfseries What About the Independents Then:\par}
\vspace{0.4cm}
{\Large Investigating Small Large Language Models in the UK\par}
\vspace{1.2cm}

{\large\textbf{Phillip Charles O'Brien (Eden Eldith)}\par}
\vspace{0.3cm}
{\small Independent AI researcher and model developer, Gosport, United Kingdom\par}
\vspace{0.2cm}
{\small ORCID: \href{https://orcid.org/0009-0007-3961-1182}{0009-0007-3961-1182}\par}

\vfill
{\small 31 October 2026 $\cdot$ Version 1.1.1\par}
\vspace{1cm}
\end{titlepage}
\hypersetup{pageanchor=true}
\pagenumbering{arabic}

\begin{abstract}
The United Kingdom is investing in sovereign AI, yet a map of AI firms
is not necessarily a map of the people who can build it. I follow
a fixed Hugging Face search for English-language, PyTorch text-generation
base models below one billion parameters: \CensusModels{} repositories from
\CensusAccounts{} publishing accounts on 30 October 2026. From saved model and
profile evidence, I identify \UKBuilders{} UK-based builders across
\UKBuiltModels{} model entries, including \ScratchBuilders{} builders of
\ScratchModels{} from-scratch entries. \UKInstitutional{} of the \UKBuilders{} builders have
university or institutional routes. \textbf{I am the one identified
strict-independent builder, having developed WiggleGPT in Gosport}, with no university
attendance or institutional affiliation. The complete account appendix
retains all institutional, overseas, derivative, untrained and unresolved
cases. Hugging Face, GitHub and Zenodo provide a practical route to discover
independent domestic capability that company-centred inventories leave
outside their populations. Finding, supporting and retaining those builders
is part of developing sovereign AI.
\end{abstract}

\textbf{Keywords:} independent research; UK AI capability; small language models; model discovery; sovereign AI.

\section{Capability that existing inventories may not see}
British AI sovereignty depends on identifying the people who can already
design, train and develop models in the United Kingdom. An independent
researcher who builds a language model from randomly initialised parameters
has created domestic technical capability, whether the work was performed
in an academic laboratory, a funded company or a home workshop. If national
mapping begins with registered businesses and university networks, such
individuals need not appear in the initial population at all.

The government's \emph{Artificial Intelligence Sector Study 2025}, published
on 29 October 2026, measures commercial AI activity through UK-registered
companies. Its methods already use frontier language models and agentic
research tools for discovering and classifying firms. Its stated population
excludes non-registered individual contractors \cite{sectorstudy}.
\textbf{The absence of an independent developer from a company census is
a consequence of its population definition, not evidence that the developer
does not exist.} That boundary matters when an economic inventory becomes
the starting point for understanding national technical capability.

The Council for Science and Technology's advice, published on 19 October
2026, calls for a Sovereignty Plan informed by a comprehensive view of
the domestic AI ecosystem \cite{cst}. Will that map search Hugging Face,
GitHub, Zenodo and individual builders, or reproduce the institutions it
already knows? The practical question is who has already built domestic
models outside university and institutional routes, and how to find them.

The author trained WiggleGPT, approximately 124 million parameters, from
scratch in Gosport on personally owned consumer GPUs \cite{wigglegpt,wigglecard}.
Six earlier manual and AI-assisted investigations did not identify another
matching independent UK builder. The author does not believe he is the
only one: this search is built to find the others. The purpose is to discover potential collaborators
and people whose demonstrated skills warrant research backing.

\section{A bounded, evidence-led discovery method}
The 30 October 2026 snapshot contains \textbf{917 model repositories from
495 publishing accounts}. The saved query selected Text Generation,
English, PyTorch, the displayed under-1B parameter range, Base only, and
Recently Updated, with no name search. The ten API pages and 31 website
pages agreed on membership. That membership remains frozen; later live
searches do not replace it.

The PyTorch filter is intentional. In the author's 31 October exploration,
adding Transformers under the final combination reduced the visible result
to 552 and hid WiggleGPT; removing that additional library constraint
restored it. Removing PyTorch in a later configuration produced 7,022
results. These are recorded live observations illustrating filter sensitivity,
not alternative census populations or estimates of all UK models.

Classification begins with the downloaded CSV/JSONL records, model cards and
profiles. Clear organisations, derivatives and overseas connections can be
recorded immediately. Existing investigations are reused; targeted additional
sources are needed only for unanswered fields that could change a relevant
classification. Provenance distinguishes new-initialisation training from
fine-tuning, continued pretraining, distilled or reused weights, mirrors and
untrained checkpoints. A catalogue's Base tag is not proof of original
training. Saved checkpoint comparisons inform the GPT-2 and T5 derivative
cases; no repeat weight downloads were needed for this revision.

\emph{UK-built} concerns the builder's location while developing and running
the model. Overseas rented GPUs are recorded separately. Dated study or
employment overlapping development is distinguished from direct run-location
evidence and from a current profile with no historical link. \emph{University
ever} includes incomplete attendance anywhere. \emph{Institutional route}
includes university, laboratory, company or relevant industry affiliation,
including a technology day job when the model was a home project.
Employment does not itself demonstrate that the employer paid for training.
Affiliation at build time, project backing, funding and compute source remain
separate fields.

The strict-independent group requires at least one genuinely from-scratch
natural-language model built in the UK, no university attendance, and no
institutional affiliation or backing, with independence affirmatively
supported. A missing biography is not proof of independence. Counts use
models for provenance, accounts for the complete inventory, and deduplicated
people for researcher routes. Exclusion reasons can overlap; everyone remains
visible in the full 495-account companion appendix.

\section{Results}
\subsection{Complete cohort and account coverage}
The website and API identifier sets matched exactly: 917 model repositories,
with no duplicate identifiers or missing pagination results. The Hub
classified 409 of their 495 publisher accounts as users and 86 as
organisations, accounting for 748 and 169 model repositories respectively.
These platform account types are not findings about funding independence.
Recorded parameter counts ranged from 2 to 998,008,576; the very small
counts are retained as catalogue observations rather than interpreted as
complete language-model sizes.

Table~\ref{tab:coverage} gives the complete account-level outcomes.
The 245 unresolved accounts constitute 49.5\% of publishers and account
for 419 of the 917 repositories. All remain in the released table.
The mixed category includes both international collaborations and
conflicting personal-location evidence. The non-UK category also includes
clearly labelled organisational locations or professional affiliations;
it is not a count of 235 individually verified non-UK residents.

\begin{table}[htbp]
\centering
\caption{Location evidence for all publishing accounts. Model counts in
this table inherit the publisher's category; model-specific UK connections
are reported separately below.}
\label{tab:coverage}
\begin{tabular}{lrr}
\toprule
Account evidence category & Accounts & Model repositories \\
\midrule
Explicit UK public location & 8 & 8 \\
UK affiliation only & 1 & 1 \\
Mixed or conflicting locations & 6 & 30 \\
Non-UK geographic evidence & 235 & 459 \\
Unresolved & 245 & 419 \\
\midrule
Total & 495 & 917 \\
\bottomrule
\end{tabular}
\end{table}

\subsection{Eight individually published UK findings}
The eight accounts with explicit UK personal locations each contributed
one repository to the cohort. Table~\ref{tab:individuals} lists every case.
The location links point to the public identity evidence; the references
support the model descriptions. Current location is not automatically the
location of an earlier training run.

\begingroup
\small
\setlength{\tabcolsep}{5pt}
\begin{longtable}{@{}>{\raggedright\arraybackslash}p{0.26\linewidth}>{\raggedright\arraybackslash}p{0.18\linewidth}>{\raggedright\arraybackslash}p{0.49\linewidth}@{}}
\caption{All eight individual publisher accounts with explicit UK locations.}
\label{tab:individuals}\\
\toprule
Researcher and release & Public location & Documented work \\
\midrule
\endfirsthead
\toprule Researcher and release & Public location & Documented work \\
\midrule
\endhead
\bottomrule
\endfoot
Phillip C. O'Brien (Eden Eldith)\newline
\path{edeneldith/WiggleGPT}
& \href{https://garden-backend-three.vercel.app/}{Gosport}
& 123.7M parameters; pretrained from scratch on OpenWebText, with a
learnable oscillating MLP activation. One consumer GPU at a time, first
RTX 3070 then RTX 5060 Ti. Explicitly describes independent work without
institutional funding \cite{wigglecard}. \\[0.6em]
Ahmad Bader\newline\path{Ahmadbader17/tiny-gpt}
& \href{https://github.com/Ahmadbader17}{Cambridge}
& 5.85M GPT-2-style model implemented and trained from scratch on a laptop
CPU. TinyStories, 77.8M training tokens, 7.9 hours on an Intel Core Ultra 9
185H. Learning/demo use; no project-funding statement \cite{tinygpt}. \\[0.6em]
Sam Schorb\newline\path{iimaginary/schorbGPT-medium}
& \href{https://github.com/sam-schorb}{Sheffield}
& GPT-2 architecture and tokenizer, described as a web-text base model
without fine-tuning. Initialisation history, hardware and funding are
not specified in the reviewed sources \cite{schorb}. \\[0.6em]
Sam Dobson\newline\path{sdobson/tinystories-llama-15m}
& \href{https://github.com/samdobson}{UK}
& 15M Llama-style TinyStories model, with an implementation adapted from
\texttt{llama2.c} and 100,000 reported iterations. Initialisation and
actual hardware are not explicitly stated \cite{dobson}. \\[0.6em]
Sripaad Srinivasan\newline\path{Isotonic/bullet-points-generator}
& \href{https://github.com/sripaad}{London}
& T5 conditional bullet-point generator. The autogenerated card claims
from-scratch training on 35,000 source/target rows; its initialisation
history is not otherwise established \cite{isotonic}. \\[0.6em]
Rajesh Kanna V\newline\path{XBOT-RK/distilgpt2-wiki-qa}
& \href{https://github.com/vrajeshtrichy}{UK}
& Explicitly fine-tuned from DistilGPT2 for Wiki-QA. The Base-only
catalogue selection did not exclude this derivative \cite{xbot}. \\[0.6em]
Sally Sonia Simmons\newline\path{SallySims/equibert-generator}
& \href{https://github.com/SallySims}{UK}
& T5-base fine-tuned for conditional text generation. Separate encoder
architecture descriptions elsewhere in the card are not attributed to
this generator \cite{sally}. \\[0.6em]
Abdullatif K\"oksal\newline\path{akoksal/LongForm-OPT-125M}
& \href{https://akoksal.com/}{London}
& OPT instruction fine-tuning using LongForm data, produced through
academic collaboration. The publisher currently works at Google DeepMind
in London; this does not establish the earlier training site \cite{longform}. \\
\end{longtable}
\endgroup

Ahmad Bader's release and WiggleGPT provide direct, documented examples
of training from scratch on consumer hardware. They contribute different
kinds of evidence. The former is a small standard-architecture educational
model trained on a laptop CPU. WiggleGPT is an architecture experiment
with a learned oscillating activation, a 600,000-step pretraining record,
released checkpoints and a stated absence of institutional funding.
Its recorded pretraining ran from 23 October to 12 November 2025, with
later release updates and task adaptations. This dated research record
is substantive capability evidence even though its creator has no
registered AI company.

The remaining personal releases demonstrate why implementation,
pretraining, fine-tuning, architecture and funding should be separate
fields. They contribute to the discovery record without being forced
into a single category of ``independent from-scratch LLM''.

\subsection{UK connections within collective releases}
Following named researchers beyond publisher addresses added eight
repositories with a person's public UK location. These comprise three
Outpost Bio Waypoint models through coauthor Neythen Treloar's London
location; three PolyPythias 160M seed releases through joint first author
Pietro Lesci's public Cambridge profiles; BLOOMZ-560M through coauthor
Samuel Albanie in London; and BLOOM-560M-intermediate through dataset
and applications contributor Daniel van Strien's UK profile
\cite{waypoint,polypythias,bloomz,bloom}.

The resulting 16 repositories span 11 publisher accounts and identify
12 named people, counting the author's public names as one person.
Fifteen repositories have an identified publisher/developer or model-paper
coauthor with UK location evidence. The sixteenth is retained explicitly
as a dataset/applications contribution, not evidence that van Strien
personally trained that checkpoint. Coauthorship does not establish who
operated the training run.

The Waypoint releases model microbiome taxonomic sequences, not ordinary
English prose. Their English metadata tag admitted them to the original
query; the task-specific evidence identifies the difference. The remaining
13 of these 16 repositories concern natural-language generation, including
task-specific and multilingual releases. Collective or company research
is retained without describing it as unaffiliated solo work.

Thirteen further repositories have UK institutional-affiliation evidence:
Augusto Corr\^ea's Oxford-linked SmolLM fine-tune; GPT-Neo-125M through
named author Sid Black's UK institutional role; and eleven original
Pythia-suite releases through Herbie Bradley's Cambridge affiliation
\cite{correa,gptneo,pythia}. Bradley's documented contribution was WinoBias
analysis, and his current personal site locates him in San Francisco.
These rows are kept in an affiliation view rather than assigned current
UK residence. Seven other reviewed collective releases had no established
model-specific UK link, despite UK participation elsewhere in the larger
organisation or model family.

\subsection{A location conflict and partial training record}
\path{AiRukua/BabyKilat} describes a custom hybrid model trained from
scratch, released at 6,000 of 50,000 planned steps \cite{babykilat}.
Its author's linked GitHub profile says Southampton, whereas LinkedIn
identifies Ambon, Indonesia and records Southampton study. The conflicting
location evidence is retained. The card gives estimated training times
for several GPUs, which do not identify the device actually used.
This case is visible in the broader 30-repository UK-connected view,
with both the location conflict and partial-training status explicit.

\subsection{Catalogue filters and substantive evidence}
Thirty-seven model-level assessments cover the personal findings,
collective cases and unresolved related variants. Twenty-three of all
917 catalogue matches contain a nonempty declared \texttt{base\_model}
field. Directly inspected cards include fine-tuned checkpoints despite
the Base-only filter. Some model cards describe a different task from
the catalogue task; others are untrained or debugging artefacts.

Parameter metadata also requires interpretation. For example,
\path{ml-resources/MobileMoE-S-QAT} appears in the under-one-billion
cohort, while its card describes 272M active parameters and 1.3B total
parameters \cite{mobilemoe}. This is a property of the retrieved catalogue
record, not a reason to silently rewrite the original cohort. The release
retains every original match and puts the substantive distinctions in
separate, source-supported annotations.


\section{Discoverable models, undiscovered builders}
WiggleGPT is publicly available with weights, code and an archived paper.
The \emph{State of AI Report} presents itself as a widely read and trusted
independent annual overview of AI progress \cite{stateofaiabout}.
Its 2026 edition, dated 8 October and covering the preceding
12 months, contains no mention of WiggleGPT or its author in its searchable
text \cite{stateofai2026}. The earliest saved Zenodo record gives an original
publication date of 25 November 2025 and a deposit created on 8 December
2025 \cite{wiggleoriginal}. The work therefore predates the report and falls
within its stated reporting window. The omission concerns the already
released original work.

The author's search experience makes the mechanism tangible. Generic
Gosport AI searches often surfaced institutional pages or jobs instead of
his work. Yahoo surfaced his Grokipedia entry where Google did not; he reads
this as relevant results being downranked. In his earlier agent-assisted
searches, broad questions about independent British model builders produced
answers that they did not exist, supported by familiar official archives.
Agents given this Hugging Face search found WiggleGPT immediately.
Asking universities who does AI also
privileges the people already connected to those universities.

The implication is direct: a capability map that consults only company
registers and established institutions cannot establish the absence of
independent capability. It has not searched the population in which that
capability may reside. Released checkpoints, training code and archival
research provide a practical alternative entry point. They can reveal
skills before the builder gains an employer, grant or incorporated company.

\section{Sovereign AI is also a talent-discovery problem}
The Sovereign AI Fund describes a \pounds500 million programme, with equity
investments typically between \pounds1 million and \pounds10 million, alongside
compute, research, procurement and talent support for companies
\cite{sovereignfaq}. Its visa scheme reimburses eligible costs when supported
UK companies recruit international specialists into UK research roles
\cite{sovereignoffer}. Its investment criteria also allow companies to
commit to establishing a substantial, lasting UK presence, rather than
requiring all recipients to be headquartered here already
\cite{sovereignfaq}. That support exists for attracting capability from abroad.
What parallel mechanism actively identifies comparable
technical talent already working independently in Britain, before company
formation? Recruitment and domestic discovery are complementary national
capabilities; a national strategy should provide an active route for both.

International research partnerships can strengthen domestic capability
through the exchange of architectural
ideas, training methods and technical improvements, with British builders
contributing in return. Talent recruitment is one route; reciprocal
technical collaboration is another. Discovering independent capability
already present in Britain gives those partnerships more people and
expertise to work with.

Girolami and Wooldridge's 2023 sovereignty editorial distinguishes domestic
from-scratch development from adaptation and licensing. Its large-scale
scenario envisages roughly 100 researchers; licensing alone does not confer
meaningful ownership \cite{girolami2023}. The Turing Institute's June 2026
analysis again proposes a medium-scale sovereign base model with control
over data, training and weights \cite{turingsovereignty}. Its Project t0 is a
Qwen2.5-based fine-tuning project \cite{turingt0}. The Institute therefore
already uses the model-repository ecosystem in which independent builders
publish their own work.

I have already developed, trained, documented and publicly released
WiggleGPT in Britain as a working sovereign base model. The project
demonstrates the capacity to design,
train and publish a model at a much smaller independently financed scale,
using less compute over more time. Scaled using national compute, it could provide the starting
point for, or replace, the proposed
Turing sovereign base model.

My historical argument is that WiggleGPT is ``the first truly sovereign AI
model made in Britain''. In six earlier investigations, I found no other
publicly documented case meeting the combination of domestic original
development, independence, from-scratch training and public release.

Those searches and
the original release chronology provide the source chain for the priority
argument; Girolami and Wooldridge provide its sovereignty framework.

Its task-adapted results are 90.1\% on SST-2 public validation and 59.5\% on
RACE-middle test \cite{wigglecard}. GPT-1 reported 91.3\% and 62.9\%,
respectively \cite{radford2018}. The historical comparison describes different
evaluation splits and training coverage: WiggleGPT's RACE adaptation used
middle-school data, while GPT-1 used both middle and high school. These
results demonstrate substantial task learning. The underlying model-building
skill is already present.

The scale of the underlying training is also worth recognising. My raw
OpenWebText corpus was approximately 22.5 GB. The GPT-2 paper reports
40 GB of cleaned WebText \cite{radford2019}. Substantial language-model
training is already possible in a home workshop. National compute would
extend an existing research capability.

A useful model does not need to lead every benchmark to have commercial
value. Performance close enough for a customer's task can be worth paying
for, especially when the model is cheaper to train and host. Small models
therefore deserve assessment on the tasks they perform, their operating
costs and the control they give their users. Funding decisions that assume
useful model development must begin at the scale of a frontier laboratory
miss this route to practical value.

An early example that informed my thinking was a DeepSeek report I
recorded in February 2025: roughly US\$400 of API use to obtain outputs for
fine-tuning a model. The cost
component in that example is access to generated training material. Access
to existing models, data and methods can make a further stage of development
affordable to people with much smaller budgets. This is a concrete reason
for international partnerships to share technical resources with domestic
builders: the exchange can enable work that would otherwise be beyond
their reach.

\subsection{Recursive self-improvement and the cost of delay}
\label{sec:rsi}
Recursive self-improvement (RSI) begins with a model helping to build the
next model. Useful coding capability already allows AI to contribute to
implementation, training, evaluation and further development under human
direction. Improvements can then feed back into the tools used for the
following round of development. The extreme case, in which improvement
outpaces human understanding, is a later possibility in that progression.
The practical development loop is already available to builders today.

The progression from AI through RSI towards artificial general intelligence
(AGI) and artificial superintelligence (ASI) is central to the strategic
argument. A country able to participate in successive rounds of development
can build expertise, produce useful systems and contribute to the next
advance. A country that remains a customer depends on the terms set by the
producers. Donald Trump expressed the competitive stakes directly in a
September 2026 call with Jensen Huang: ``whoever wins AI, wins''
\cite{trump2026ai}. This is why finding and supporting domestic builders is
urgent: their work contributes to a process in which capability can compound.

Delay has a cost. While an independent builder waits for recognition,
research opportunities, collaborators and investment can move elsewhere.
The next development then takes place around the people and infrastructure
that received support. Waiting until builders acquire conventional
credentials or form a company can mean discovering their capability only
after the opportunity to develop it domestically has passed. Britain's
starting population must include the people doing the work now, and its
support must arrive in time to help them do the next part.

The consequences extend to national security. More capable AI can change
the balance between attack and defence, including the effectiveness of
weapons and countermeasures on which present deterrence depends. Decisions
about domestic model-building therefore affect future freedom of action
and bargaining power. Britain needs the capacity to develop and understand
the technology as it advances.

The National Security Strategy 2025 seeks a baseline sovereign AI capability
\cite{nss2025}. Sensitive systems depend on retained technical knowledge,
continuity of skilled staff and control over external dependencies. Public
investment should secure long-term domestic commitment through contracts
and retention arrangements, with explicit responsibilities for continuity
and protection of sensitive knowledge. This is a concrete way to keep the
expertise available over five- and ten-year horizons, as engineering
organisations already seek to do. Sovereignty policy should account for the
capacity that remains available after the initial investment announcement.

\subsection{Sovereignty has to survive procurement and operation}
The NHS provides a concrete test. In 2023 NHS England awarded the Federated
Data Platform and Associated Services contract to a consortium led by
Palantir Technologies. The government's account assigns control of the
data to the NHS and describes Palantir as processing it under NHS
instructions \cite{nhspalantir2025}. Operational sovereignty also requires
the ability to maintain, change and replace the systems through which
those instructions are carried out. What can Britain operate independently
if access to an external supplier changes? Which domestic people and teams
can build the replacement? Those are capability questions that procurement
and national AI policy need to answer together.

Maintenance is part of the same problem. The National Audit Office's
investigation of WannaCry recorded earlier warnings about obsolete
software, critical patch alerts in March and April 2017, and no formal
mechanism before the attack for checking whether local NHS organisations
had acted. Infected organisations had unpatched or unsupported Windows
systems \cite{naowannacry}. Published advice had not become dependable
operational practice. Sovereignty requires people able to keep systems
working and secure throughout their life, with the resources to do so.

The same applies to finding vulnerabilities. Describing known threats
does not establish the ability to discover new weaknesses in our own
systems, test them and repair them. Model-building competence, active
security research and control of the deployment environment belong in
the same capability map. Sensitive health and defence systems need that
practical capacity behind claims of British control.

Access restrictions elsewhere demonstrate the strategic importance attached
to this technology. US advanced-chip controls explicitly address national
security \cite{bis2024,bis2026}; Anthropic's dated 2026 account describes
China, Russia and Iran as unsupported countries \cite{anthropicregions2026}.
Why are overseas access boundaries treated as strategically important while
independently built domestic capabilities remain invisible to recognition
systems? The question concerns identifiable discovery and support decisions.

The local political context reinforces it. Rupert Lowe called for sovereign
AI on 24 October 2026, having separately commented on Gosport on 5 October
\cite{loweai,lowegosport}. A town receiving national attention already
contained my publicly released original model. I want an effective
discovery route without having to become a political campaign's example.


\section{From inventory to opportunity}
The constructive purpose is to make talent visible. A public index of
independent builders would let research groups, compute providers, companies
and public bodies find documented work, identify collaborators and determine
where support could accelerate it. Technical capability should be
discoverable before a researcher learns to market a company or builds a
large social-media audience.

Educational access helps explain why university networks are an incomplete
starting point. The Department for Education's July 2026 release follows
English state-school pupils aged 15 in 2020/21 to age 19 in 2024/25. It reports
higher-education progression of 44.7\% overall and 28.6\% among
pupils eligible for free school meals; the corresponding rate for White
British boys eligible for free school meals is 13.0\% \cite{dfe2026}.
The June 2026 \emph{Independent Inquiry into White Working Class Educational
Outcomes} also identifies working-class disadvantage beyond FSM eligibility
\cite{wwci}. University networks therefore leave a substantial route into
technical learning and capability outside their normal field of view.

Learning to train a small language model can start much earlier than a
research degree. Students can initialise weights, observe loss decline,
generate imperfect text, change the architecture and see the result. A
character-level Shakespeare model gives this process a concrete form.
Schools, after-school clubs, community programmes and apprenticeships could
teach Python, PyTorch, training, evaluation and deployment. Kaggle provides
free, quota-limited GPU notebooks \cite{kagglegpu}; Google's September 2026 announcement
extends premium Colab benefits to Google AI subscribers \cite{colab2026}.
Publicly explaining these routes would help people use access they may
already have.

The author's route combined practical hardware experience, growing up
helping in his father's computer shop, with self-directed learning; he dates
his start in Python to 17 January 2025. His work was assisted by AI. He
directed Astra to carry out WiggleGPT's task-adaptation work with the
instruction to make it do what GPT-1 did. The resulting task evaluations are
documented in the model release \cite{wigglecard}. AI assistance can put
capability in places established discovery systems have not thought to look.

Regional compute access, machine-learning apprenticeships and technical hubs
would connect those beginnings to sustained research. Gosport illustrates
why the search should extend beyond the familiar university and technology
centres. Wider infrastructure planning must account for electricity
distribution, cooling and long-term resilience. Funding only established
performers misses the next capable researcher. Regional access and practical
learning routes create a path for that researcher to emerge and be recognised.

Publication access is another part of discoverability. arXiv requires
endorsement for a first submission to a new endorsement category, with
automatic endorsement possible in specified circumstances; institutional
email can assist the process \cite{arxivendorsement}. Zenodo and model
repositories offer additional routes for independent work to become public.
A national inventory should follow reproducible outputs across these venues
instead of treating participation in one academic publication route as the
boundary of research capability.

\section{Evidence coverage}
The census covers the specified 917-model search on 30 October 2026.
Public evidence leaves 151 account locations and 385 model origins
unresolved. Six of the seven UK development entries use dated professional
or study-location inferences; WiggleGPT has direct project-location evidence.
Lesci's joint-first-author role establishes participation in PolyPythias;
the exact division of training operations is unrecorded. Project sponsorship
and some hardware details remain unknown. The appendix retains the sources,
evidence levels and open fields for every account.

\section{What about the independents?}

Britain's ability to develop sovereign AI depends on whether it knows who
can already build the technology. The independent case found here is public
and inspectable. The people paid to discover sovereign capability should
have found publicly released independent work on the main model-hosting
platform. That failure is \emph{a major indictment of their ability to
capture sovereign AI talent and capability}. The criticism is constructive: the aim is for the
people responsible for discovery to do better, and for the builders they
find to have a route to collaboration and sustained support.

\textbf{If Britain wants to know what sovereign AI capability it already
has, it must look not only where companies and universities report it, but
where independent individuals actually build and publish it.} Becoming a
producer rather than merely a consumer of AI includes identifying, teaching,
supporting and retaining those individuals.

\section*{Availability and authorship}
Code, paper and companion appendix:
\url{https://github.com/Eden-Eldith/uk-small-model-census}.
Frozen metadata, source-linked classifications and reproducible counts:
\url{https://huggingface.co/datasets/edeneldith/uk-small-model-census}.
Version 1.1.1 includes the corrected manuscript and matching Zenodo package
for author deposit.

Phillip Charles O'Brien supplied the research question, filters, interpretation,
authorial narrative and publication decisions. Astra and Codex (OpenAI) and
Claude (Anthropic) assisted with research and drafting under his direction.
The author created WiggleGPT, the confirmed strict-independent case.
Private recordings, transcripts and research working notes are not released.
\printbibliography
\end{document}
